% segment-local-id: OR03-S011
% segment-id: d90.orig.v1.tr03.seg0011
% license: CC BY-SA 4.0
% provenance: Materi asesmen asli berbahasa Indonesia, disusun secara independen.

\section{Ujian akhir kumulatif}
\label{orig03:final}

% stable-id: d90.orig.v1.tr03.final.exercise.0001
\begin{exercise}[Eksistensi komposit dengan regularisasi kuadratik]
\label{orig03:final:exercise:0001}
Misalkan $h:\mathbb R^m\to\mathbb R$ konveks, kontinu, dan tak negatif;
$r:\mathbb R^n\to(-\infty,+\infty]$ proper, semikontinu-bawah, konveks, dan
tak negatif; $A:\mathbb R^n\to\mathbb R^m$ linear; serta $\lambda>0$.
Buktikan bahwa
\[
 F(x)=h(Ax)+r(x)+\frac\lambda2\lVert x\rVert^2
\]
mempunyai tepat satu peminimum. Nyatakan bagian argumen yang membuktikan
keberadaan dan bagian yang membuktikan ketunggalan.
\end{exercise}

% stable-id: d90.orig.v1.tr03.final.hint.0001
\paragraph{Petunjuk bertahap.}
\label{orig03:final:hint:0001}
\textbf{Tahap 1.} Gunakan ketaknegatifan dua suku pertama untuk membatasi
himpunan sublevel.
\textbf{Tahap 2.} Periksa apa yang terjadi pada modulus konveksitas kuat ketika
fungsi konveks ditambahkan pada kuadrat $\lambda\lVert x\rVert^2/2$.

% stable-id: d90.orig.v1.tr03.final.answer.0001
\paragraph{Jawaban singkat.}
\label{orig03:final:answer:0001}
$F$ proper, semikontinu-bawah, dan
$F(x)\ge\lambda\lVert x\rVert^2/2$, sehingga ia koersif dan mencapai minimum.
Selain itu, $F$ bersifat $\lambda$-konveks kuat, sehingga peminimumnya tunggal.

% stable-id: d90.orig.v1.tr03.final.solution.0001
\paragraph{Solusi lengkap.}
\label{orig03:final:solution:0001}
Karena $r$ proper, terdapat $x_0$ dengan $r(x_0)<\infty$; fungsi $h$ bernilai
hingga, jadi $F(x_0)<\infty$ dan $F$ proper. Komposisi $h\circ A$ kontinu dan
konveks, sementara $r$ semikontinu-bawah. Maka $F$ semikontinu-bawah. Dari
ketaknegatifan,
\[
 F(x)\ge\frac\lambda2\lVert x\rVert^2\longrightarrow+\infty
 \quad\text{ketika }\lVert x\rVert\to\infty.
\]
Jadi sebuah barisan peminimum berada dalam himpunan terbatas; dalam dimensi
hingga kita dapat mengambil subsekuens konvergen, dan semikontinuitas-bawah
menunjukkan bahwa limitnya mencapai infimum. Ini membuktikan keberadaan.

Fungsi kuadrat adalah $\lambda$-konveks kuat. Penambahan dua fungsi konveks
mempertahankan modulus itu, jadi $F$ juga $\lambda$-konveks kuat. Jika ada dua
peminimum berbeda, ketaksamaan konveksitas kuat pada titik tengah akan memberi
nilai yang lebih kecil secara ketat daripada nilai minimum, kontradiksi. Maka
peminimum tunggal.

% stable-id: d90.orig.v1.tr03.final.exercise.0002
\begin{exercise}[Pencarian balik untuk model proksimal]
\label{orig03:final:exercise:0002}
Misalkan $f$ konveks dengan gradien $L$-Lipschitz, $r$ proper
semikontinu-bawah dan konveks, serta $x\in\operatorname{dom}r$. Untuk $M>0$,
definisikan
\[
 x_M=\arg\min_y\left\{
 r(y)+f(x)+\langle\nabla f(x),y-x\rangle
 +\frac M2\lVert y-x\rVert^2\right\}.
\]
Algoritme menggandakan $M$ sampai uji majorisasi
\[
 f(x_M)\le f(x)+\langle\nabla f(x),x_M-x\rangle
 +\frac M2\lVert x_M-x\rVert^2                         \tag{A}
\]
terpenuhi. Buktikan bahwa pencarian berhenti secara hingga, lalu buktikan
sertifikat penurunan
\[
 (f+r)(x_M)\le(f+r)(x)-\frac M2\lVert x_M-x\rVert^2.
\]
Nyatakan residu gradien proksimal yang dapat dihitung dari langkah ini.
\end{exercise}

% stable-id: d90.orig.v1.tr03.final.hint.0002
\paragraph{Petunjuk bertahap.}
\label{orig03:final:hint:0002}
\textbf{Tahap 1.} Lemma penurunan menjamin (A) untuk semua $M\ge L$.
\textbf{Tahap 2.} Model yang diminimalkan bersifat $M$-konveks kuat dan
nilainya di $y=x$ sama dengan $(f+r)(x)$.

% stable-id: d90.orig.v1.tr03.final.answer.0002
\paragraph{Jawaban singkat.}
\label{orig03:final:answer:0002}
Penggandaan akhirnya menghasilkan $M\ge L$, sehingga (A) berlaku. Konveksitas
kuat model memberi nilai model di $x_M$ paling besar
$(f+r)(x)-M\lVert x_M-x\rVert^2/2$; (A) memindahkan batas itu ke objektif
sebenarnya. Residu yang lazim adalah
$G_M(x)=M(x-x_M)$.

% stable-id: d90.orig.v1.tr03.final.solution.0002
\paragraph{Solusi lengkap.}
\label{orig03:final:solution:0002}
Lemma penurunan menyatakan untuk semua $y$,
\[
 f(y)\le f(x)+\langle\nabla f(x),y-x\rangle
 +\frac L2\lVert y-x\rVert^2.
\]
Jadi begitu urutan penggandaan mencapai $M\ge L$, uji (A) pasti diterima;
jumlah penggandaan hingga.

Tuliskan model di dalam kurung sebagai $Q_M(y)$. Ia $M$-konveks kuat dan
$x_M$ adalah peminimumnya. Dengan memakai $0\in\partial Q_M(x_M)$ dalam
ketaksamaan konveksitas kuat,
\[
 Q_M(x)\ge Q_M(x_M)+\frac M2\lVert x-x_M\rVert^2.
\]
Karena $Q_M(x)=(f+r)(x)$ dan uji (A) memberi
$(f+r)(x_M)\le Q_M(x_M)$, kedua ketaksamaan menghasilkan sertifikat yang
diminta. Vektor $G_M(x)=M(x-x_M)$ dapat dihitung dari langkah. Ia nol tepat
bila $x=x_M$, yang melalui kondisi optimalitas ekuivalen dengan
$0\in\nabla f(x)+\partial r(x)$.

% stable-id: d90.orig.v1.tr03.final.exercise.0003
\begin{exercise}[KKT untuk regularisasi analisis dan kendala norma]
\label{orig03:final:exercise:0003}
Untuk $\tau>0$ dan $R>0$, pertimbangkan
\[
 \min_x\ \frac12\lVert Ax-b\rVert^2+\tau\lVert x\rVert_1
 \quad\text{dengan syarat}\quad \lVert Bx-d\rVert_2\le R.
\]
Anggap terdapat $\bar x$ dengan $\lVert B\bar x-d\rVert_2<R$. Turunkan sistem
KKT lengkap. Tuliskan subgradien norma Euclidean baik ketika
$Bx-d\ne0$ maupun ketika $Bx-d=0$, dan jelaskan mengapa sistem KKT cukup untuk
optimalitas global.
\end{exercise}

% stable-id: d90.orig.v1.tr03.final.hint.0003
\paragraph{Petunjuk bertahap.}
\label{orig03:final:hint:0003}
\textbf{Tahap 1.} Gunakan pengali $\nu\ge0$ untuk fungsi kendala
$q(x)=\lVert Bx-d\rVert_2-R$.
\textbf{Tahap 2.} Terapkan aturan rantai subdiferensial pada norma dan pisahkan
kasus argumen nol.

% stable-id: d90.orig.v1.tr03.final.answer.0003
\paragraph{Jawaban singkat.}
\label{orig03:final:answer:0003}
Terdapat $s\in\partial\lVert x\rVert_1$, $q\in\partial\lVert Bx-d\rVert_2$,
dan $\nu\ge0$ sehingga
\[
 0=A^\top(Ax-b)+\tau s+\nu B^\top q,
\]
\[
 \lVert Bx-d\rVert_2\le R,qquad
 \nu(\lVert Bx-d\rVert_2-R)=0.
\]
Jika $z\ne0$, $\partial\lVert z\rVert_2=\{z/\lVert z\rVert_2\}$; jika
$z=0$, subdiferensialnya adalah bola satuan. Slater membuat KKT perlu dan,
karena masalah konveks, juga cukup.

% stable-id: d90.orig.v1.tr03.final.solution.0003
\paragraph{Solusi lengkap.}
\label{orig03:final:solution:0003}
Lagrangian adalah
\[
 \mathcal L(x,\nu)=\tfrac12\lVert Ax-b\rVert^2
 +\tau\lVert x\rVert_1+\nu(\lVert Bx-d\rVert_2-R).
\]
Kondisi Slater yang diberikan menjamin keberadaan pengali optimal dan aturan
KKT tanpa celah dual. Stasioneritas adalah inklusi
\[
 0\in A^\top(Ax-b)+\tau\partial\lVert x\rVert_1
 +\nu B^\top\partial\lVert Bx-d\rVert_2.
\]
Memilih $s$ dan $q$ dari kedua subdiferensial mengubahnya menjadi persamaan
pada jawaban singkat. Secara koordinat,
$s_i=\operatorname{sign}(x_i)$ jika $x_i\ne0$ dan $s_i\in[-1,1]$ jika
$x_i=0$. Untuk norma Euclidean, $q=z/\lVert z\rVert$ saat $z=Bx-d\ne0$,
sedangkan sembarang $q$ dengan $\lVert q\rVert\le1$ sah saat $z=0$.

Sistem dilengkapi oleh kelayakan primal, kelayakan dual $\nu\ge0$, dan
kelonggaran komplementer. Objektif dan fungsi kendala konveks; maka setiap
pasangan yang memenuhi sistem membentuk sertifikat titik pelana Lagrangian
dan memberi peminimum global.

% stable-id: d90.orig.v1.tr03.final.exercise.0004
\begin{exercise}[Distribusi kepentingan dengan momen kedua minimum]
\label{orig03:final:exercise:0004}
Misalkan $g_1,\ldots,g_N$ adalah vektor tak nol yang diketahui. Untuk peluang
$p_i>0$, $\sum_i p_i=1$, ambil $I\sim p$ dan
$G=g_I/(Np_I)$. Di antara semua distribusi $p$, tentukan distribusi yang
meminimalkan $\mathbb E\lVert G\rVert^2$, hitung nilai minimumnya, dan tulis
varians minimum terhadap mean $\bar g=N^{-1}\sum_i g_i$.
\end{exercise}

% stable-id: d90.orig.v1.tr03.final.hint.0004
\paragraph{Petunjuk bertahap.}
\label{orig03:final:hint:0004}
\textbf{Tahap 1.} Masalah skalar yang harus diminimalkan adalah
$\sum_i\lVert g_i\rVert^2/p_i$ pada simpleks.
\textbf{Tahap 2.} Gunakan pengali Lagrange atau Cauchy--Schwarz, lalu kurangi
$\lVert\bar g\rVert^2$ dari momen kedua.

% stable-id: d90.orig.v1.tr03.final.answer.0004
\paragraph{Jawaban singkat.}
\label{orig03:final:answer:0004}
\[
 p_i^\star=\frac{\lVert g_i\rVert}{\sum_j\lVert g_j\rVert},
 \qquad
 \min_p\mathbb E\lVert G\rVert^2
 =\frac1{N^2}\left(\sum_i\lVert g_i\rVert\right)^2.
\]
Varians minimumnya adalah nilai terakhir dikurangi $\lVert\bar g\rVert^2$.

% stable-id: d90.orig.v1.tr03.final.solution.0004
\paragraph{Solusi lengkap.}
\label{orig03:final:solution:0004}
Ketakbiasan mengikuti dari
$\mathbb E G=\sum_i p_i g_i/(Np_i)=\bar g$. Momen keduanya ialah
\[
 \mathbb E\lVert G\rVert^2
 =\frac1{N^2}\sum_i\frac{\lVert g_i\rVert^2}{p_i}.
\]
Cauchy--Schwarz memberi
\[
 \left(\sum_i\lVert g_i\rVert\right)^2
 =\left(\sum_i\frac{\lVert g_i\rVert}{\sqrt{p_i}}\sqrt{p_i}\right)^2
 \le\left(\sum_i\frac{\lVert g_i\rVert^2}{p_i}\right)
      \left(\sum_i p_i\right).
\]
Kesetaraan terjadi tepat ketika
$\lVert g_i\rVert/\sqrt{p_i}$ sebanding dengan $\sqrt{p_i}$, yakni
$p_i\propto\lVert g_i\rVert$. Karena semua vektor tak nol, distribusi itu
strictly positif dan layak. Substitusi memberi nilai minimum. Identitas
$\mathbb E\lVert G-\bar g\rVert^2=
\mathbb E\lVert G\rVert^2-\lVert\bar g\rVert^2$ memberi varians yang diminta.

% stable-id: d90.orig.v1.tr03.final.exercise.0005
\begin{exercise}[Laju rata-rata descent cermin stokastik]
\label{orig03:final:exercise:0005}
Misalkan $C$ konveks, $h$ terdiferensialkan dan $1$-konveks kuat terhadap
norma $\lVert\cdot\rVert$, serta $f$ konveks pada $C$. Dengan filtrasi
$\mathcal F_k$, anggap
\[
 \mathbb E[G_k\mid\mathcal F_k]=g_k\in\partial f(x_k),
 \qquad
 \mathbb E[\lVert G_k\rVert_*^2\mid\mathcal F_k]\le G^2.
\]
Iterasi diberikan oleh
\[
 x_{k+1}\in\arg\min_{x\in C}
 \{\eta\langle G_k,x\rangle+D_h(x,x_k)\}.
\]
Jika $x^\star$ meminimalkan $f$ dan
$D_h(x^\star,x_0)\le D^2/2$, buktikan untuk
$\bar x_K=K^{-1}\sum_{k=0}^{K-1}x_k$ dan
$\eta=D/(G\sqrt K)$ bahwa
\[
 \mathbb E[f(\bar x_K)-f(x^\star)]\le\frac{DG}{\sqrt K}.
\]
\end{exercise}

% stable-id: d90.orig.v1.tr03.final.hint.0005
\paragraph{Petunjuk bertahap.}
\label{orig03:final:hint:0005}
\textbf{Tahap 1.} Dari ketaksamaan tiga titik dan konveksitas kuat $h$,
turunkan batas satu langkah dengan suku $\eta^2\lVert G_k\rVert_*^2/2$.
\textbf{Tahap 2.} Ambil ekspektasi bersyarat, jumlahkan, gunakan Jensen pada
$f(\bar x_K)$, lalu optimalkan langkah konstan.

% stable-id: d90.orig.v1.tr03.final.answer.0005
\paragraph{Jawaban singkat.}
\label{orig03:final:answer:0005}
Ketaksamaan dasar memberi
\[
 \eta\,\mathbb E[f(x_k)-f(x^\star)]
 \le\mathbb E[D_h(x^\star,x_k)-D_h(x^\star,x_{k+1})]
 +\frac{\eta^2G^2}{2}.
\]
Setelah dijumlahkan dan dibagi $\eta K$, batasnya
$D^2/(2\eta K)+\eta G^2/2$; pilihan langkah yang diberikan membuatnya
$DG/\sqrt K$.

% stable-id: d90.orig.v1.tr03.final.solution.0005
\paragraph{Solusi lengkap.}
\label{orig03:final:solution:0005}
Optimalitas langkah cermin dan identitas tiga titik menghasilkan, untuk
$u\in C$,
\[
 \eta\langle G_k,x_{k+1}-u\rangle
 \le D_h(u,x_k)-D_h(u,x_{k+1})-D_h(x_{k+1},x_k).
\]
Tambahkan $\eta\langle G_k,x_k-x_{k+1}\rangle$ pada kedua ruas. Karena
$D_h(x_{k+1},x_k)\ge\lVert x_{k+1}-x_k\rVert^2/2$, Young memberi
\[
 \eta\langle G_k,x_k-u\rangle
 \le D_h(u,x_k)-D_h(u,x_{k+1})
 +\frac{\eta^2}{2}\lVert G_k\rVert_*^2.
\]
Ambil $u=x^\star$ dan ekspektasi bersyarat. Karena $x_k$ terukur terhadap
$\mathcal F_k$, ketakbiasan memberi
$\mathbb E[\langle G_k,x_k-x^\star\rangle\mid\mathcal F_k]
=\langle g_k,x_k-x^\star\rangle\ge f(x_k)-f(x^\star)$.
Maka ketaksamaan pada jawaban singkat berlaku.

Menjumlahkan untuk $k=0,\ldots,K-1$ meneleskopkan divergensi dan memberi
\[
 \frac1K\sum_{k=0}^{K-1}\mathbb E[f(x_k)-f(x^\star)]
 \le\frac{D^2}{2\eta K}+\frac{\eta G^2}{2}.
\]
Konveksitas dan Jensen membatasi ruas kiri dari bawah oleh
$\mathbb E[f(\bar x_K)-f(x^\star)]$. Substitusi
$\eta=D/(G\sqrt K)$ membuat kedua suku kanan sama dengan
$DG/(2\sqrt K)$, sehingga hasil mengikuti.

% stable-id: d90.orig.v1.tr03.final.exercise.0006
\begin{exercise}[Kontraksi resolven operator monoton kuat]
\label{orig03:final:exercise:0006}
Misalkan $A$ maksimal monoton dan $\mu$-monoton kuat dengan $\mu>0$.
Buktikan bahwa untuk setiap $\gamma>0$,
\[
 \lVert J_{\gamma A}x-J_{\gamma A}y\rVert
 \le\frac1{1+\gamma\mu}\lVert x-y\rVert.
\]
Simpulkan bahwa $A$ mempunyai paling banyak satu nol dan, jika nol
$x^\star$ ada, iterasi $x_{k+1}=J_{\gamma A}x_k$ konvergen linear kepadanya.
\end{exercise}

% stable-id: d90.orig.v1.tr03.final.hint.0006
\paragraph{Petunjuk bertahap.}
\label{orig03:final:hint:0006}
\textbf{Tahap 1.} Tetapkan $p=J_{\gamma A}x$, $q=J_{\gamma A}y$, lalu tulis
anggota $Ap$ dan $Aq$ dari definisi resolven.
\textbf{Tahap 2.} Terapkan monotoni kuat dan Cauchy--Schwarz; kemudian ambil
$y=x^\star$.

% stable-id: d90.orig.v1.tr03.final.answer.0006
\paragraph{Jawaban singkat.}
\label{orig03:final:answer:0006}
Monotoni kuat memberi
\[
 \langle p-q,x-y\rangle\ge(1+\gamma\mu)\lVert p-q\rVert^2,
\]
yang bersama Cauchy--Schwarz menghasilkan kontraksi. Karena
$J_{\gamma A}x^\star=x^\star$,
\[
 \lVert x_k-x^\star\rVert
 \le(1+\gamma\mu)^{-k}\lVert x_0-x^\star\rVert.
\]

% stable-id: d90.orig.v1.tr03.final.solution.0006
\paragraph{Solusi lengkap.}
\label{orig03:final:solution:0006}
Dari definisi resolven,
\[
 \frac{x-p}{\gamma}\in Ap,
 \qquad
 \frac{y-q}{\gamma}\in Aq.
\]
Monotoni kuat memberikan
\[
 \left\langle p-q,\frac{x-p}{\gamma}-\frac{y-q}{\gamma}\right\rangle
 \ge\mu\lVert p-q\rVert^2.
\]
Setelah dikalikan $\gamma$ dan ditata ulang, ini menjadi ketaksamaan pada
jawaban singkat. Jika $p\ne q$, Cauchy--Schwarz dan pembagian oleh
$\lVert p-q\rVert$ memberi batas Lipschitz; bila $p=q$ batasnya trivial.

Jika $u$ dan $v$ keduanya nol $A$, monotoni kuat memberi
$0\ge\mu\lVert u-v\rVert^2$, sehingga $u=v$. Untuk nol yang ada,
$J_{\gamma A}x^\star=x^\star$. Penerapan berulang faktor kontraksi
$q=(1+\gamma\mu)^{-1}<1$ memberi batas linear yang dinyatakan.

% stable-id: d90.orig.v1.tr03.final.exercise.0007
\begin{exercise}[Sinkhorn sebagai maksimisasi blok dual]
\label{orig03:final:exercise:0007}
Misalkan $a_i,b_j>0$ dan total massanya sama. Untuk $\varepsilon>0$, tinjau
\[
 D(f,g)=a^\top f+b^\top g
 -\varepsilon\sum_{i,j}
 \exp\!\left(\frac{f_i+g_j-C_{ij}}{\varepsilon}\right).
\]
Buktikan bahwa, dengan $g$ tetap, pemaksimal unik terhadap setiap $f_i$ adalah
\[
 f_i=\varepsilon\log a_i-\varepsilon
 \log\sum_j\exp\!\left(\frac{g_j-C_{ij}}{\varepsilon}\right),
\]
dan turunkan pembaruan analog untuk $g_j$. Simpulkan bahwa penskalaan Sinkhorn
adalah maksimisasi koordinat blok eksak dan nilai dual tidak menurun pada
setiap setengah langkah. Jelaskan ketidaktunggalan bersama akibat gauge.
\end{exercise}

% stable-id: d90.orig.v1.tr03.final.hint.0007
\paragraph{Petunjuk bertahap.}
\label{orig03:final:hint:0007}
\textbf{Tahap 1.} Nolkan turunan parsial terhadap satu $f_i$ dan periksa tanda
turunan keduanya.
\textbf{Tahap 2.} Ulangi untuk $g_j$ dan amati transformasi
$(f,g)\mapsto(f+c\1,g-c\1)$.

% stable-id: d90.orig.v1.tr03.final.answer.0007
\paragraph{Jawaban singkat.}
\label{orig03:final:answer:0007}
Persamaan stasioner blok $f$ ialah
\[
 a_i=\sum_j\exp((f_i+g_j-C_{ij})/\varepsilon),
\]
yang memberi rumus pada soal; rumus blok $g$ diperoleh dengan menukar baris
dan kolom. Setiap turunan kedua blok negatif, jadi pembaruan adalah maksimum
unik dan tidak menurunkan $D$. Pasangan bersama hanya unik hingga penambahan
konstanta pada $f$ dan pengurangan konstanta yang sama pada $g$.

% stable-id: d90.orig.v1.tr03.final.solution.0007
\paragraph{Solusi lengkap.}
\label{orig03:final:solution:0007}
Untuk $g$ tetap,
\[
 \frac{\partial D}{\partial f_i}
 =a_i-\sum_j\exp((f_i+g_j-C_{ij})/\varepsilon).
\]
Menetapkannya nol dan mengeluarkan faktor $\exp(f_i/\varepsilon)$ memberi
rumus eksplisit pada soal. Turunan kedua adalah
\[
 -\frac1\varepsilon\sum_j
 \exp((f_i+g_j-C_{ij})/\varepsilon)<0,
\]
dan tidak ada suku silang antara $f_i$ yang berlainan ketika $g$ tetap. Jadi
rumus itu adalah pemaksimal blok unik. Secara simetris,
\[
 g_j=\varepsilon\log b_j-\varepsilon
 \log\sum_i\exp\!\left(\frac{f_i-C_{ij}}{\varepsilon}\right).
\]
Mengganti satu blok dengan pemaksimal eksaknya tidak dapat menurunkan nilai
fungsi dual. Dalam variabel skala eksponensial, kedua rumus adalah pembaruan
baris dan kolom Sinkhorn.

Jika $f$ ditambah $c\1$ dan $g$ dikurangi $c\1$, semua jumlah
$f_i+g_j$ tetap. Suku linear juga tetap karena total massa $a$ dan $b$ sama.
Jadi nilai dual dan kopling tidak berubah; inilah kebebasan gauge yang membuat
pemaksimal pasangan tidak tunggal sebelum normalisasi.

% stable-id: d90.orig.v1.tr03.final.exercise.0008
\begin{exercise}[Pemilihan pemisahan untuk model stokastik terstruktur]
\label{orig03:final:exercise:0008}
Pertimbangkan
\[
 \min_{x\in C}\quad
 s(x)+\tau\lVert Wx\rVert_1,
 \qquad
 s(x)=\mathbb E_\xi[\ell_\xi(A_\xi x-b_\xi)],
\]
dengan $C=\{x:\lVert x\rVert_2\le R\}$, $\tau>0$, dan $s$ konveks
terdiferensialkan. Tuliskan inklusi optimalitas primal dan sistem primal--dual
dengan variabel $y$ untuk suku $\ell^1$. Bandingkan kapan gradien proksimal
stokastik langsung layak dipakai dan kapan pemisahan primal--dual lebih
sesuai. Berikan sepasang residu KKT dan asumsi orakel minimum yang diperlukan
untuk klaim konvergensi stokastik.
\end{exercise}

% stable-id: d90.orig.v1.tr03.final.hint.0008
\paragraph{Petunjuk bertahap.}
\label{orig03:final:hint:0008}
\textbf{Tahap 1.} Gabungkan $\partial(\tau\lVert W\cdot\rVert_1)$ dengan
kerucut normal $N_C$, lalu gunakan aturan rantai melalui $W^\top$.
\textbf{Tahap 2.} Bandingkan proksimal gabungan dengan proyeksi ke $C$ dan
proksimal konjugat norma $\ell^1$ yang dapat dihitung secara terpisah.

% stable-id: d90.orig.v1.tr03.final.answer.0008
\paragraph{Jawaban singkat.}
\label{orig03:final:answer:0008}
Optimalitas ekuivalen dengan keberadaan $y$ sehingga
\[
 0\in\nabla s(x)+W^\top y+N_C(x),
 \qquad y\in\tau\partial\lVert Wx\rVert_1.
\]
Gradien proksimal langsung sesuai bila
$\operatorname{prox}_{\eta(\tau\lVert W\cdot\rVert_1+\delta_C)}$ mudah;
pemisahan primal--dual sesuai bila hanya proyeksi ke $C$ dan proyeksi dual ke
bola $\ell^\infty$ yang mudah. Residu dapat berupa jarak kedua inklusi ke nol.
Orakel harus tidak bias secara bersyarat, bermomen kedua terkendali, dan
langkah harus memenuhi syarat penjumlahan yang sesuai.

% stable-id: d90.orig.v1.tr03.final.solution.0008
\paragraph{Solusi lengkap.}
\label{orig03:final:solution:0008}
Dengan aturan rantai konveks,
\[
 \partial\bigl(\tau\lVert W\cdot\rVert_1\bigr)(x)
 =W^\top\bigl(\tau\partial\lVert Wx\rVert_1\bigr).
\]
Karena kendala ditulis sebagai indikator $\delta_C$, syarat Fermat adalah
\[
 0\in\nabla s(x)+W^\top
       \bigl(\tau\partial\lVert Wx\rVert_1\bigr)+N_C(x).
\]
Memperkenalkan $y\in\tau\partial\lVert Wx\rVert_1$ menghasilkan sistem pada
jawaban singkat. Secara ekuivalen,
$Wx\in\partial(\tau\lVert\cdot\rVert_1)^*(y)$, dan fungsi konjugat itu adalah
indikator bola $\{y:\lVert y\rVert_\infty\le\tau\}$.

Jika proksimal gabungan
$\tau\lVert W\cdot\rVert_1+\delta_C$ tersedia dalam bentuk tertutup atau dapat
dihitung murah, langkah gradien proksimal stokastik memisahkan bagian halus
$s$ secara langsung. Untuk $W$ umum, proksimal gabungan biasanya tidak
terurai menjadi penyusutan lunak lalu proyeksi. Metode primal--dual menjaga
$Wx$ sebagai operator linear dan hanya membutuhkan proyeksi ke bola $C$ serta
proyeksi $y$ ke bola $\ell^\infty$ berjari-jari $\tau$; karena itu ia lebih
sesuai pada keadaan tersebut.

Satu pasangan residu matematis adalah
\[
 r_{\mathrm p}(x,y)=
 \operatorname{dist}\bigl(0,\nabla s(x)+W^\top y+N_C(x)\bigr),
\]
\[
 r_{\mathrm d}(x,y)=
 \operatorname{dist}\bigl(Wx,\partial(\tau\lVert\cdot\rVert_1)^*(y)\bigr),
\]
dengan nilai tak hingga bila $y$ berada di luar domain konjugat. Keduanya nol
tepat pada pasangan KKT\@. Untuk analisis stokastik, penduga gradien
$G_k$ sekurang-kurangnya harus memenuhi
$\mathbb E[G_k\mid\mathcal F_k]=\nabla s(x_k)$ dan batas momen kedua atau
varians yang sesuai. Pada skema Robbins--Monro dasar, langkah positif lazimnya
memenuhi $\sum_k\eta_k=\infty$ dan $\sum_k\eta_k^2<\infty$; klaim laju dengan
langkah tetap memerlukan asumsi tambahan seperti reduksi varians atau menerima
lantai galat. Norma operator $W$ juga harus masuk ke batas langkah
primal--dual, bukan diabaikan.
